% =============================================================================
% Conclusiones
% =============================================================================
\section{Conclusiones}

% Síntesis de hallazgos
\subsection{Síntesis de hallazgos}

El aprendizaje profundo ha pasado de ser experimental a una herramienta con evidencia sólida de eficacia clínica. Las CNN pueden igualar o superar el rendimiento humano en diagnóstico radiológico, alcanzando un AUC del 93.2\% en clasificación de gliomas y del 89.6\% en detección de cáncer de mama. U-Net se ha consolidado como estándar para segmentación y el aprendizaje por transferencia ha permitido superar la escasez de datos médicos anotados. Más de 1,000 dispositivos basados en IA han sido autorizados por la FDA, el 88.2\% en radiología.

% Valoración crítica del equipo
\subsection{Valoración crítica del equipo}

Pese a la alta precisión, la implementación masiva enfrenta desafíos críticos. El problema de la \enquote{Caja Negra} persiste como barrera para la confianza clínica, por lo que consideramos imperativo adoptar la IA como \enquote{Inteligencia Aumentada}: un copiloto que complementa al especialista sin reemplazar su autonomía diagnóstica.

Además, muchos modelos heredan sesgos algorítmicos por falta de diversidad en los datos, lo que podría perpetuar desigualdades en poblaciones marginadas. La \enquote{Deriva de la IA}, donde el rendimiento decae por cambios en equipos o protocolos, exige marcos de gobernanza con monitoreo continuo y recalibración periódica.

% Perspectivas futuras del área
\subsection{Perspectivas futuras del área}

Las fronteras de investigación se centran en las siguientes tendencias:

\begin{enumerate}
    \item \textbf{IA Explicable}: Sistemas que permitan comprender qué regiones de la imagen influyeron en el diagnóstico, eliminando el efecto de caja negra.
    \item \textbf{Modelos Fundacionales}: Transición hacia modelos generales entrenados en conjuntos multimodales (rayos X, MRI y CT) para diagnósticos transversales.
    \item \textbf{NLP}: Integración de notas clínicas textuales con datos de imagen para generar reportes automatizados.
    \item \textbf{Aprendizaje Federado}: Entrenamiento en redes de múltiples hospitales sin que los datos privados salgan de su institución de origen.
    \item \textbf{AR y VR}: Mejora de la visualización durante cirugías mínimamente invasivas y planeamiento preoperatorio.
\end{enumerate}

En conclusión, la IA en radiología se encuentra en fase de expansión donde la colaboración entre radiólogos, desarrolladores y reguladores será decisiva para una implementación equitativa, segura y efectiva.
